\SourceLeaf{038}
\pdfbookmark[0]{Ein weiterer neuer Beweis}{second-paper}
\thispagestyle{empty}
\begin{center}
\vspace*{.25in}
{\large\bfseries EIN WEITERER NEUER BEWEIS}\par
\vspace{.18in}
DES SATZES\par
\vspace{.35in}
{\Large\bfseries DASS JEDES POLYNOM\par
\medskip }\par
\vspace{.25in}
IN EINER VARIABLEN\par
\vspace{.25in}
{\bfseries IN REELLE FAKTOREN ERSTEN ODER ZWEITEN GRADES}\par
\vspace{.25in}
ZERLEGT WERDEN KANN\par
\vspace{.6in}
V O N\par
\vspace{.3in}
{\large CARL FRIEDRICH GAUSS}\par
\vspace{.35in}
DER KÖNIGLICHEN GESELLSCHAFT DER WISSENSCHAFTEN VORGELEGT, 1815, DEZEMBER 7.\par
\vspace{.65in}
\rule{.88\linewidth}{.3pt}\par
\medskip
{\small Commentationes societatis regiae scientiarum Gottingensis recentiores. Vol.~\textsc{iii}.\par}
\medskip
Göttingen \textsc{mdcccxvi}.\par
\medskip
\rule{.88\linewidth}{.3pt}
\end{center}
\clearpage
